\section{Introduction}

Active inference treats perception and action as coupled inference under a generative model. An agent maintains beliefs about hidden causes of observations, evaluates policies by expected free energy, and selects actions by a softmax over policy values weighted by policy precision \todocite{active inference and expected free energy}. In this framework, policy precision is not a decorative temperature parameter. It determines whether the agent acts decisively on the currently preferred policy or preserves a more exploratory distribution over plausible policies.

Affective active-inference accounts give policy precision psychological content. In Hesp et al.'s account of deeply felt affect, affective valence is tied to expected action precision and subjective model fitness: positive affect corresponds to an expectation that the current generative model is working well enough to support confident action, while negative affect corresponds to model failure or reduced expected precision \todocite{Hesp et al. Deeply Felt Affect}. This converts affect from an externally added reward bonus into a metacognitive signal about whether the agent can trust its own model.

Social interaction makes this model-fitness problem naturally multi-factorial. A focal agent interacting with several partners may have a reliable model of one partner, an unreliable model of another, and a confidently adversarial model of a third. A single global affective precision signal collapses these distinctions. In social environments, the relevant question is often not ``Do I feel good?'' or ``Has this partner rewarded me?'' but ``How well does my model of this partner currently work?'' A predictable exploiter can have high model fitness even if interacting with them is unrewarding; a volatile ally can have low model fitness even if their average payoff is favorable.

This paper tests a partner-specific extension of affective precision in a repeated trust task. The central claim is that affective precision gates deployment of otherwise available social knowledge. It is not modeled as a partner liking score, a reward cache, or a literal trust scalar. It is an auxiliary partner-local estimate of the reliability of the focal agent's own social model. The signal updates from prediction error about the partner's action and modulates the policy posterior used to choose actions or partners.

The project contributes four elements. First, it factorizes Hesp-style expected action precision over social partners, making affect a set of partner-local model-fitness estimates rather than a single global action-model state. Second, it implements the mechanism around official \texttt{inferactively-pymdp==1.0.0} agents: state and policy inference remain in standard partner-local POMDP agents, while beta tracking is an external task-owned module. Third, it tests the mechanism with readouts that separate model fitness from reward, deployment from belief accuracy, social choice from payoff, and stress amplification from recovery success. Fourth, it reports a productive boundary condition: abrupt betrayal can turn the same precision channel into misdeployment by sharpening action under a wrong or insufficiently recalibrated post-switch model.

The resulting claim is deliberately narrower than ``affect improves payoff.'' In open graded regimes, partner-specific precision can lower entropy and improve payoff. In reliability-versus-reward tests, it follows predictive reliability more strongly than payoff. In social-choice tests, it changes partner-selection distributions before payoff is informative. Under abrupt betrayal, however, affective precision can lower entropy while worsening payoff. This combination is theoretically useful because it shows that the mechanism is active but not uniformly beneficial.
